\section{Some Group Theory}
\label{app:gen:groups}

\begin{definition}
    A \emph{\green{group}} $(G,\cdot)$ consists of a (possibly infinite) set of
    elements and a binary op $\cdot : G \times G \to G$ such that
    \begin{enumerate}
        \item $(a \cdot b) \cdot c = a \cdot (b \cdot c)$
        \item There exists an $\id_G \in G$ s.t. $\id_G \cdot\;a = a \cdot \id_G = a$
        \item Each $a \in G$ is invertible with $a \cdot a^{-1} = a^{-1} a =
        \id_G$
    \end{enumerate}
\end{definition}

\begin{definition}
    A group $(G,\cdot)$ is called \emph{\green{abelian}}, if for all $a,b \in G$
    it holds that $a \cdot b = b \cdot a$, its contrapose is called a
    \emph{\green{non-abelian}} group.
\end{definition}

Recall that we say a group is \emph{generated} by a set of elements if we can
write any other element as the concatenation of the generating set. Here by
example:

\begin{definition}
    The \emph{\green{symmetric group}} of $n$ points $S_n$ is given by
    \begin{equation}
        S_n 
        = 
        \left\langle 
        s_1, \ldots, s_{n-1} 
        \;\middle\vert\;    
        \begin{aligned}
            s_i^2 &= id \\
            s_i s_{i+1} s_i &= s_{i+1} s_i s_{i+1} \\
            s_i s_j &= s_j s_i, \vert i - j \vert \geq 2
        \end{aligned}
        \;\right\rangle.
    \end{equation}
\end{definition}

\begin{definition}[Representation theory of the symmetric group]
    For $n \geq 2$ the symmetric group $S_n$ is equipped with two
    one-dimensional representations, the \emph{trivial representation} where
    every generator is equal to the identity, and the \emph{sign representation}
    which equips each generator with its relative permutation parity.
\end{definition}

\begin{definition}
    The \emph{\green{General Linear Group}} $\GL{n}$ of degree $n$ is the set of
    $n\times n$ matrices that are invertible, and its operation $\cdot$ is that
    of matrix multiplication. 
\end{definition}

\section{Some Algebraic Topology}
\label{app_gen:homotopy}

Let spaces $X,Y$ be topological in the following, and let all functions be
continuous unless stated otherwise.

\begin{definition}
    A \emph{\green{homotopy}} between two continuous functions $f,g : X \to Y$
    is a continuous function
    \begin{equation}
        H : X \times [0,1] \to Y \text{ such that } H(x,0) = f(x) \text{ and } H(x,1) = g(x)
    \end{equation}
\end{definition}

\begin{definition}
    A \emph{\green{loop}} with basepoint $x_0$ on $X$ is a function
    \begin{equation}
        \gamma : [0,1] \to X \quad \text{such that} \gamma(0) = \gamma(1) = x_0
    \end{equation}
\end{definition}

\section{Tensor Decompositions \& Algorithms}

\begin{definition}
    The \emph{\green{singular value decomposition}} of a complex operator $A$ of shape $(m,n)$ is a decomposition of the form
    \[
        A = U\Sigma V^\ast
    \]
    where
    \[
        U \in \U{m \times m}, \quad \Sigma \in \Mat{m }{n}{\R} \quad \text{and} \quad V \in \U{n \times n}
    \]
    and $\Sigma$ is further diagonal, in fact is it made up of the \emph{singular values} of $A$.

    This also enables us to perform a low rank appoximation of an operator by truncating $\Sigma$ to only keep its $r$ largest singular values. 
\end{definition}